\documentclass[12pt,a4paper]{article}
\usepackage[utf8]{inputenc}
\usepackage[T1]{fontenc}
\usepackage[polish]{polski}
\usepackage{amsmath}
\usepackage{amsfonts}
\usepackage{amssymb}
\usepackage{booktabs}
\usepackage{geometry}
\usepackage{enumitem}
\usepackage{float}
\usepackage{listings}
\usepackage{xcolor}
\usepackage{array}
\usepackage{multirow}

\geometry{margin=2cm}

\definecolor{codebg}{rgb}{0.97,0.97,0.97}

\lstset{
    basicstyle=\ttfamily\small,
    backgroundcolor=\color{codebg},
    frame=single,
    framesep=4pt,
    showstringspaces=false,
}

\title{Systemy czasu rzeczywistego, Projekt 4 — Tło teoretyczne}
\author{Piotr Gugnowski, wydział EiTI, nr indeksu 320690}
\date{Czerwiec 2026}

\pagestyle{empty}

\begin{document}

\maketitle

% ━━━━━━━━━━━━━━━━━━━━━━━━━━━━━━━━━━━━━━━━━━━━━━━━━━━━━━━━━━━━━━━━━━━━━━━━━━━━
\section*{Dodatek A — Tło teoretyczne}
% ━━━━━━━━━━━━━━━━━━━━━━━━━━━━━━━━━━━━━━━━━━━━━━━━━━━━━━━━━━━━━━━━━━━━━━━━━━━━

% ─────────────────────────────────────────────────────────────────────────────
\subsection*{A.0 Słowniczek — co znaczą podstawowe pojęcia}

\textbf{Magistrala} (\textit{bus}) — wspólna para przewodów łącząca wszystkie
urządzenia w~sieci. Każdy sygnał jest widoczny przez wszystkich — jak tablica
ogłoszeń, z~której każdy może czytać.

\medskip
\textbf{Węzeł} (\textit{node}) — jedno urządzenie podłączone do magistrali
(czujnik, sterownik, komputer).

\medskip
\textbf{Ramka} (\textit{frame}) — jedna paczka danych. Poza właściwymi danymi
zawiera nagłówek (numer wiadomości, ile jest danych), sumę kontrolną i~inne pola.
Jak koperta pocztowa: prócz listu jest adres, znaczek itp.

\medskip
\textbf{Stany recessive i~dominant.}
Magistrala CAN ma w~każdej chwili jeden z~dwóch stanów:
\begin{itemize}
    \item \textit{Recessive} (z~łac.\ cofający się) = logiczne~1 — stan
          \emph{bierny}: żadne urządzenie nic nie robi, magistrala
          „odpoczywa" na napięciu $\approx$2,5\,V na obu przewodach.
    \item \textit{Dominant} (z~łac.\ panujący) = logiczne~0 — stan
          \emph{aktywny}: co najmniej jedno urządzenie aktywnie steruje
          napięciami na przewodach.
\end{itemize}
\textbf{Zasada nadrzędna:} jeden dominant bierze górę nad wszystkimi
recessive. Jeśli urządzenie A wysyła dominant (0) jednocześnie gdy B wysyła
recessive (1), na magistrali widać 0. Jak głośny krzyk zagłuszający szept.

\medskip
\textbf{Sygnałowanie różnicowe — skąd nazwa?}
Zwykły kabel cyfrowy przesyła napięcie jednego przewodu względem masy:
0\,V to zero, 3,3\,V to jedynka. Problem: każdy impuls elektromagnetyczny
(od silnika, telefonu, błyskawicy) bezpośrednio zmienia to napięcie.

CAN używa \emph{dwóch} przewodów: CAN\_H (High) i~CAN\_L (Low).
Informacja jest zakodowana w~\textbf{różnicy} napięć między nimi:

\begin{center}
\begin{tabular}{lrrr}
\toprule
\textbf{Stan} & CAN\_H & CAN\_L & Różnica H$-$L \\
\midrule
\textit{recessive} (1) & $\approx 2{,}5$\,V & $\approx 2{,}5$\,V & $\approx 0$\,V \\
\textit{dominant}  (0) & $\approx 3{,}5$\,V & $\approx 1{,}5$\,V & $\approx 2$\,V \\
\bottomrule
\end{tabular}
\end{center}

Zakłócenie uderza jednakowo w~oba przewody biegnące obok siebie — podnosi
napięcie na obu o~tę samą wartość, więc różnica H$-$L pozostaje niezmieniona.
Stąd nazwa: \textbf{różnicowe} — liczymy różnicę, nie wartość bezwzględną.

\medskip
\textbf{Ciągnięcie napięcia.}
W~stanie recessive nikt nic nie robi: rezystory zamontowane na końcach kabla
(rezystory terminacyjne) utrzymują oba przewody pasywnie na $\approx$2,5\,V.
W~stanie dominant układ scalony (transceiver — nadajnik/odbiornik CAN)
włącza wewnętrzne tranzystory, które:
\begin{itemize}
    \item \textbf{ciągną CAN\_H w~górę} — ku napięciu zasilania ($\approx$3,5\,V),
    \item \textbf{ciągną CAN\_L w~dół} — ku masie ($\approx$1,5\,V).
\end{itemize}
„Ciągnięcie" to standardowy termin elektroniczny: tranzystor musi dostarczyć
prąd, żeby wymusić napięcie wbrew rezystorom — jak ciągnięcie liny wbrew
sprężynie.

\medskip
\textbf{Arbitracja — kto nadaje pierwszy?}
CAN to sieć bez żadnego nadrzędnego koordynatora: każde urządzenie może zacząć
nadawać w~dowolnej chwili. Gdy dwa węzły zaczną jednocześnie — kolizja.

CAN rozwiązuje ją automatycznie: każde nadające urządzenie \emph{słucha}
magistrali podczas własnej transmisji i~porównuje to co wysyła z~tym co widzi.
Numer ID wiadomości jest wysyłany bit po bicie od najstarszego.
Jeśli urządzenie wysłało 1 (recessive), ale na magistrali widzi 0 (dominant) —
znaczy to, że inne urządzenie wysyła tam 0, czyli ma niższy numer ID.
Urządzenie z~wyższym ID milknie i~czeka. Wygrywa ten z~niższym numerem ID
— i~nie traci przy tym żadnych danych, bo dominanty wygrywają bez zniszczenia
ramki. Stąd: \textbf{niższy numer ID = wyższy priorytet}.

\medskip
\textbf{Bit stuffing (wstawianie bitów synchronizujących).}
Odbiornik nie ma osobnego kabla zegarowego. Synchronizuje swój wewnętrzny zegar
obserwując \emph{zmiany} stanu na magistrali — każde przejście $0\to1$ lub
$1\to0$ jest jak uderzenie metronomu. Problem: przy długim ciągu identycznych
bitów (np.\ 10 jedynek z~rzędu) przez cały ten czas nie ma żadnej zmiany
i~odbiornik może „zgubić rytm".

Rozwiązanie: nadajnik automatycznie \textbf{wstawia jeden odwrócony bit po
każdych 5~bitach tej samej wartości}. Odbiornik zna tę zasadę i~usuwa te
dodatkowe bity. Takie bity pomocnicze to \textbf{bity stuffujące}.

\medskip
\textbf{CRC} (\textit{Cyclic Redundancy Check} — cykliczna suma kontrolna).
Matematyczna metoda wykrywania przekłamań. Nadajnik oblicza 15-bitową liczbę
kontrolną ze wszystkich bitów wiadomości i~dołącza ją na końcu. Odbiornik
oblicza tę samą liczbę ze \emph{odebranych} bitów i~porównuje — jeśli różne,
gdzieś nastąpił błąd. Szczegóły w~A.5.

\medskip
\textbf{Skróty pól ramki}: SOF, EOF, IFS, ACK, DLC, SRR, IDE, RTR —
każde wyjaśnione przy tabeli w~A.2.

% ─────────────────────────────────────────────────────────────────────────────
\subsection*{A.1 Sieć CAN — podstawy}

CAN (\textit{Controller Area Network}) to standard komunikacyjny opracowany przez
firmę Bosch w~1986~roku, pierwotnie na potrzeby motoryzacji. Dziś stosowany jest
również w~automatyce przemysłowej, medycynie i~robotyce.

\textbf{Kluczowe właściwości:}
\begin{itemize}
    \item \textbf{Każdy może nadawać w~dowolnej chwili} — nie ma żadnego
          centralnego koordynatora. Konflikty przy jednoczesnym nadawaniu
          rozstrzyga arbitracja (patrz A.0): wygrywa wiadomość z~niższym
          numerem ID, bez straty danych.
    \item \textbf{Odporność na zakłócenia elektromagnetyczne} — dzięki
          sygnałowaniu różnicowemu (dwa przewody CAN\_H i~CAN\_L;
          szczegóły w~A.0).
    \item \textbf{Dwa stany magistrali:} \textit{dominant} (logiczne~0,
          aktywnie sterowane napięcia) i~\textit{recessive} (logiczne~1,
          stan bierny). Jeden dominant zawsze wygrywa nad recessive
          (szczegóły w~A.0).
    \item Prędkość: do 1\,Mb/s (klasyczny CAN), do 8\,Mb/s (CAN~FD).
\end{itemize}

\textbf{Priorytety:} numer ID wiadomości decyduje o~priorytecie.
Niższy numer ID = więcej zer (dominantów) na początku identyfikatora =
wygrywa arbitrację. Wiadomości krytyczne (np.\ czujnik poduszki powietrznej)
dostają niskie numery ID i~zawsze trafiają na magistralę jako pierwsze.

% ─────────────────────────────────────────────────────────────────────────────
\subsection*{A.2 Struktura rozszerzonej ramki CAN (CAN~2.0B)}

CAN~2.0 definiuje dwa formaty ramek — różnią się długością numeru ID:
\begin{itemize}
    \item \textbf{Standardowy} (CAN~2.0A) — 11-bitowy numer ID
          (do 2032 różnych typów wiadomości),
    \item \textbf{Rozszerzony} (CAN~2.0B) — 29-bitowy numer ID
          (ponad 532~miliony typów wiadomości).
\end{itemize}

Zadanie 4.3 dotyczy ramki rozszerzonej z~8~bajtami danych.
Każde \textbf{pole} to wydzielony fragment ramki o~ustalonej liczbie bitów
i~ściśle określonej roli:

\begin{table}[H]
\centering
\small
\begin{tabular}{llrp{6.5cm}}
\toprule
\textbf{Pole} & \textbf{Podpole} & \textbf{Bitów} & \textbf{Co to jest} \\
\midrule
SOF & & 1 & \textit{Start of Frame} — jeden bit dominant (0) otwierający ramkę.
            Magistrala jest normalnie recessive (1); nagłe~0 budzi wszystkich
            odbiorców: „zaraz przyjdzie wiadomość". \\
\midrule
\multirow{5}{*}{Arbitration}
 & Base ID [28:18] & 11 & Starsze 11 bitów 29-bitowego numeru ID wiadomości. \\
 & SRR & 1 & \textit{Substitute Remote Request} — zawsze recessive; techniczny
             bit zgodności ze starszym formatem 11-bitowym, bez znaczenia
             merytorycznego. \\
 & IDE & 1 & \textit{Identifier Extension} — wartość recessive (1) mówi
             odbiornikowi: „to jest ramka rozszerzona (29-bitowy ID),
             nie standardowa (11-bitowy ID)". \\
 & Extension ID [17:0] & 18 & Młodsze 18 bitów 29-bitowego numeru ID. \\
 & RTR & 1 & \textit{Remote Transmission Request} — gdy ustawiony na~1,
             ramka jest \emph{prośbą} o~dane od innego węzła, nie ich
             wysłaniem. Tu zawsze~0 (wysyłamy dane, nie prosimy). \\
\midrule
\multirow{3}{*}{Control}
 & r1 & 1 & Bit zarezerwowany dla przyszłych wersji standardu; zawsze 0. \\
 & r0 & 1 & Drugi bit zarezerwowany; zawsze 0. \\
 & DLC [3:0] & 4 & \textit{Data Length Code} — liczba (0–8) mówiąca,
                    ile bajtów właściwych danych zawiera pole Data. \\
\midrule
Data & & 64 & Właściwe dane: 8~bajtów $\times$ 8~bitów = 64~bity. \\
\midrule
\multirow{2}{*}{CRC}
 & CRC sequence & 15 & 15-bitowa suma kontrolna obliczona z~wszystkich
                       wcześniejszych pól. Odbiornik liczy tę samą sumę
                       i~porównuje — niezgodność oznacza błąd transmisji. \\
 & CRC delimiter & 1 & Jeden bit recessive (1) oddzielający sumę kontrolną
                       od potwierdzenia odbioru — jak spacja między słowami. \\
\midrule
\multirow{2}{*}{ACK}
 & ACK slot & 1 & Potwierdzenie odbioru: nadajnik wysyła tu recessive (1),
                  ale \emph{każdy odbiorca, który poprawnie odebrał ramkę},
                  nadpisuje to dominant (0). Jeśli nadajnik widzi na magistrali~0
                  — wie, że ktoś odebrał wiadomość poprawnie. \\
 & ACK delimiter & 1 & Jeden bit recessive (1) kończący pole ACK. \\
\midrule
EOF & & 7 & \textit{End of Frame} — 7 bitów recessive (samych jedynek)
            oznaczających koniec ramki. \\
\midrule
IFS & & 3 & \textit{Interframe Space} — obowiązkowe 3 bity przerwy między
            ramkami, dające wszystkim węzłom czas na przetworzenie właśnie
            odebranej wiadomości zanim następna się pojawi. \\
\bottomrule
\end{tabular}
\caption{Pola rozszerzonej ramki danych CAN~2.0B przy 8~bajtach danych.}
\end{table}

\textbf{Łączna liczba bitów} dla DLC = 8 (8~bajtów danych):

\begin{align*}
L_{\text{stały}} &= \underbrace{1}_{\text{SOF}}
                 + \underbrace{(11+1+1+18+1)}_{\text{Arbitration}=32}
                 + \underbrace{(1+1+4)}_{\text{Control}=6}
                 + \underbrace{64}_{\text{Data}}
                 + \underbrace{(15+1)}_{\text{CRC}=16}
                 + \underbrace{(1+1)}_{\text{ACK}=2}
                 + \underbrace{7}_{\text{EOF}}
                 + \underbrace{3}_{\text{IFS}} \\
                 &= 1 + 32 + 6 + 64 + 16 + 2 + 7 + 3 = \mathbf{131\ \text{bitów}}
\end{align*}

Do tej liczby dochodzą jeszcze \textbf{bity stuffujące} (patrz A.3).
Bit stuffing stosuje się do regionu: SOF + Arbitration + Control + Data +
sekwencja CRC (bez separatora CRC i~dalszych pól):

\begin{equation}
N_{\text{region stuffingu}} = 1 + 32 + 6 + 64 + 15 = \mathbf{118\ \text{bitów}}
\end{equation}

% ─────────────────────────────────────────────────────────────────────────────
\subsection*{A.3 Technika bit stuffing — zasada i~zastosowanie}

\textbf{Dlaczego w~ogóle potrzebny jest bit stuffing?}

Odbiornik nie ma osobnego kabla z~sygnałem zegarowym. Synchronizuje swój
wewnętrzny zegar obserwując \emph{zmiany} stanu na magistrali — każde
przejście $0\to1$ lub $1\to0$ jest jak uderzenie metronomu. Wyobraź sobie
odczytywanie Morse'a ze słuchu: jeśli przez długi czas nie słyszysz żadnej
zmiany, gubisz rachubę. Jeśli dane zawierają 6 lub więcej identycznych bitów
z~rzędu (np.\ sześć jedynek), odbiornik może przesunąć swoje taktowanie
i~błędnie zinterpretować kolejne bity.

\textbf{Zasada:}

Po każdych \textbf{5~bitach tej samej wartości} nadajnik automatycznie
wstawia jeden bit o~\emph{przeciwnej} wartości. Ten dodatkowy bit nie niesie
żadnych danych — służy tylko do podtrzymania synchronizacji. Odbiornik zna
tę zasadę i~usuwa co szósty bit, jeśli poprzedzały go pięć identycznych.
Takie bity pomocnicze to \textbf{bity stuffujące}.

\textbf{Przykład:}
\begin{center}
\begin{tabular}{ll}
\toprule
Dane do wysłania & \texttt{1 1 1 1 1 0 0 0 0 0 1 1} \\
Na magistrali (po wstawieniu) & \texttt{1 1 1 1 1 \textbf{[0]} 0 0 0 0 0 \textbf{[1]} 1 1} \\
\bottomrule
\end{tabular}
\end{center}
Dodane bity zaznaczono nawiasami $[\,]$. Odbiornik je usuwa i~odtwarza
oryginalne dane.

\textbf{Gdzie bit stuffing obowiązuje?}
\begin{itemize}
    \item \textbf{Tak} (118~bitów łącznie): SOF, pole Arbitration, pole
          Control, pole Data, sekwencja CRC.
    \item \textbf{Nie}: separator CRC, pole ACK, EOF, IFS — te pola mają
          z~góry ustalone stałe wartości, więc odbiornik nie musi liczyć
          przejść, bo wie czego się spodziewać.
\end{itemize}

\textbf{Ile bitów stuffujących może dojść maksymalnie?}

W~najgorszym przypadku co~5 bitów wymusza wstawienie jednego stuffującego.
Dla regionu 118~bitów:

\begin{equation}
\text{Max.\ bitów stuffujących} = \left\lfloor \frac{118}{5} \right\rfloor = \mathbf{23}
\end{equation}

\textbf{Minimum} wynosi 0 — gdy dane są naprzemianne (\texttt{10101010\ldots})
i~nigdzie nie pojawia się 5~identycznych bitów z~rzędu.

% ─────────────────────────────────────────────────────────────────────────────
\subsection*{A.4 Identyfikatory CAN i~ograniczenia specyfikacji}

W~sieci CAN wiadomości nie mają adresata — każda ramka wysłana na magistralę
jest widoczna przez wszystkie węzły. Każdy węzeł sam decyduje, czy dana
wiadomość go dotyczy, na podstawie \textbf{numeru ID}. Węzeł konfiguruje
swój kontroler CAN tak, aby reagował tylko na wybrane numery ID (analogia:
stacja radiowa — każde radio słyszy wszystkie częstotliwości, ale strojysz
tylko jedną).

\textbf{Ile różnych numerów ID istnieje?}
\begin{itemize}
    \item Ramka standardowa (11~bitów na numer ID): $2^{11} = 2048$ kombinacji.
    \item Ramka rozszerzona (29~bitów na numer ID): $2^{29} = 536\,870\,912$ kombinacji.
\end{itemize}

\textbf{Ograniczenie:} identyfikatory, których \emph{pierwsze 7~bitów} to
wszystkie jedynki (1111111...), są \textbf{zarezerwowane i~zakazane}.
Powód: taka kombinacja wygląda identycznie jak specjalny sygnał błędu
(\textit{passive error flag} — 6 recessive) albo jak pole IFS (pauza między
ramkami), co mogłoby zmylić odbiorniki.

\begin{table}[H]
\centering
\begin{tabular}{lrrr}
\toprule
\textbf{Format} & \textbf{Bity ID} & \textbf{Zakazane ID} & \textbf{Dostępne ID} \\
\midrule
Standardowy & 11 & $2^{4} = 16$ & $2048 - 16 = \mathbf{2032}$ \\
Rozszerzony & 29 & $2^{22} = 4\,194\,304$ & $536\,870\,912 - 4\,194\,304 = \mathbf{532\,676\,608}$ \\
\bottomrule
\end{tabular}
\caption{Rzeczywista liczba dostępnych identyfikatorów w~sieci CAN.}
\end{table}

W~praktyce ograniczenie to jest nieistotne: typowa sieć przemysłowa używa
kilkudziesięciu do kilkuset różnych typów wiadomości — ułamek procenta
dostępnej przestrzeni.

% ─────────────────────────────────────────────────────────────────────────────
\subsection*{A.5 Suma kontrolna CRC-15/CAN — teoria i~gwarancje}

\subsubsection*{Idea: jak CRC wykrywa błędy?}

Wyobraź sobie, że wysyłasz liczbę i~chcesz sprawdzić, czy dotarła bez zmian.
Możesz z~góry obliczyć resztę z~dzielenia tej liczby przez jakąś ustaloną
stałą i~dołączyć ją jako „podpis". Odbiorca dzieli odebraną liczbę przez tę
samą stałą — jeśli reszta się zgadza, liczba prawdopodobnie jest nienaruszona.
Jeśli nie — gdzieś nastąpiła zmiana.

CRC działa dokładnie tak, ale operuje na bitach i~używa specjalnej arytmetyki
binarnej (bez przeniesień przy dodawaniu — patrz B.1). Wynikiem jest
15-bitowa reszta z~dzielenia — to właśnie suma kontrolna CRC-15.

Formalnie: ciąg bitów wiadomości traktowany jest jak wielomian (każdy bit to
współczynnik przy kolejnej potędze $x$). Nadajnik liczy resztę z~dzielenia
wielomianu wiadomości $M(x)$ przez ustalony wielomian generujący $G(x)$:
\[
R(x) = M(x) \cdot x^{15} \bmod G(x)
\]
i~dołącza resztę $R(x)$ do ramki. Odbiornik sprawdza, czy odebrany strumień
jest podzielny przez $G(x)$ bez reszty — niezerowa reszta = błąd.

\subsubsection*{Wielomian generujący CRC-15/CAN}

Wielomian generujący to matematyczna stała — rodzaj „dzielnika". Jego wybór
decyduje o~tym, jakie błędy CRC jest w~stanie wykryć. Dla CAN jest to:

\begin{equation}
G(x) = x^{15} + x^{14} + x^{10} + x^8 + x^7 + x^4 + x^3 + 1
\end{equation}

W~algorytmie z~zadania ta stała pojawia się jako liczba szesnastkowa \texttt{0x4599}
(15~bitów, bez najstarszego bitu $x^{15}$, który jest niejawny):
\[
G_{\text{hex}} = \texttt{0x4599}
= \underbrace{0100\ 0101\ 1001\ 1001}_{\text{bity 14\ldots0}}
\]

\subsubsection*{Co CRC-15/CAN gwarantuje?}

Wielomian był starannie wybrany przez Boscha pod kątem zakłóceń typowych dla
okablowania samochodowego. Gwarantuje wykrycie:
\begin{itemize}
    \item \textbf{każdego błędu 1-bitowego} — jeden przekłamany bit,
    \item \textbf{każdego błędu 2-bitowego} — dowolne dwa przekłamane bity,
    \item \textbf{każdej serii błędów} o~długości do 15~bitów z~rzędu,
    \item \textbf{każdej nieparzystej liczby przekłamanych bitów}.
\end{itemize}
Prawdopodobieństwo, że jakiś błąd \emph{nie} zostanie wykryty, wynosi
$\approx 3{,}1 \times 10^{-5}$ — czyli mniej niż 1 na 32~000 błędnych ramek
„prześlizgnie się" niezauważone.

\subsubsection*{Algorytm LFSR z~zadania}

Algorytm z~treści zadania to sprzętowa implementacja tego dzielenia, zwana
LFSR (\textit{Linear Feedback Shift Register} — liniowy rejestr przesuwający
ze sprzężeniem zwrotnym). Szczegółowy opis z~analogią taśmy produkcyjnej
znajduje się w~B.2. Zmienne algorytmu:

\begin{center}
\begin{tabular}{ll}
\toprule
\textbf{Zmienna} & \textbf{Co przechowuje} \\
\midrule
\texttt{CRC\_RG(14:0)} & 15-bitowy rejestr — bieżąca reszta dzielenia (bit 14 = najstarszy) \\
\texttt{NXTBIT} & kolejny bit wejściowy (0 lub 1) \\
\texttt{CRCNXT} & bit pomocniczy = XOR bitu wejściowego z~najstarszym bitem rejestru \\
\bottomrule
\end{tabular}
\end{center}

\begin{enumerate}
    \item \texttt{CRC\_RG = 0} (inicjalizacja — rejestr wyzerowany)
    \item Dla każdego bitu wejściowego \texttt{NXTBIT} (od SOF do końca Data):
    \begin{enumerate}[label=\alph*)]
        \item $\texttt{CRCNXT} = \texttt{NXTBIT} \oplus \texttt{CRC\_RG}(14)$
        \item przesuń rejestr w~lewo o~1:
              $\texttt{CRC\_RG}(14:1) \leftarrow \texttt{CRC\_RG}(13:0)$;
              $\texttt{CRC\_RG}(0) = 0$
        \item jeśli $\texttt{CRCNXT} = 1$: wykonaj XOR rejestru z~\texttt{0x4599}
    \end{enumerate}
    \item Wynik: zawartość \texttt{CRC\_RG} po przetworzeniu wszystkich bitów.
\end{enumerate}


% ━━━━━━━━━━━━━━━━━━━━━━━━━━━━━━━━━━━━━━━━━━━━━━━━━━━━━━━━━━━━━━━━━━━━━━━━━━━━
\section*{Dodatek B — Wymagana wiedza i~kluczowe aspekty przy implementacji}
% ━━━━━━━━━━━━━━━━━━━━━━━━━━━━━━━━━━━━━━━━━━━━━━━━━━━━━━━━━━━━━━━━━━━━━━━━━━━━

% ─────────────────────────────────────────────────────────────────────────────
\subsection*{B.1 Arytmetyka binarna bez przeniesień — dlaczego XOR jest sercem CRC}

Normalna arytmetyka: $1 + 1 = 2$ (przeniesienie do następnej cyfry).
CRC używa arytmetyki binarnej \textbf{bez przeniesień}: $1 + 1 = 0$.
To jest dokładnie operacja XOR (exclusive or — „albo jedno, albo drugie,
nie oba naraz"):
\[
0 \oplus 0 = 0,\quad 1 \oplus 0 = 1,\quad 1 \oplus 1 = 0
\]

Ta arytmetyka nosi nazwę \textbf{GF(2)} (ciało Galois rzędu~2).
Kluczowe właściwości: dodawanie i~odejmowanie to ta sama operacja (XOR),
nie ma przeniesień, nie ma ujemnych liczb. Mnożenie to AND ($1 \cdot 1 = 1$,
$1 \cdot 0 = 0$).

Ciąg bitów można interpretować jako wielomian: ciąg \texttt{1011} to
$x^3 + x + 1$ (jedynki na pozycjach 3, 1, 0). Dzielenie takiego wielomianu
przez inny w~arytmetyce GF(2) daje resztę — to właśnie CRC.

Dwie kluczowe obserwacje dla implementacji:
\begin{itemize}
    \item Przesunięcie bitów rejestru w~lewo o~1 (\texttt{<<\,1}) to mnożenie
          wielomianu przez $x$.
    \item XOR rejestru z~wielomianem generującym to odjęcie tego wielomianu
          (w~GF(2) dodawanie = odejmowanie, więc XOR wystarcza do obu operacji).
\end{itemize}
Dlatego algorytm LFSR (przesuń w~lewo + warunkowo XOR z~generatorem) to
krok po kroku implementacja dzielenia wielomianów w~GF(2).

% ─────────────────────────────────────────────────────────────────────────────
\subsection*{B.2 Rejestr LFSR — jak algorytm z~zadania działa w~praktyce}

\textbf{LFSR} (\textit{Linear Feedback Shift Register} — liniowy rejestr
przesuwający ze sprzężeniem zwrotnym) można wyobrazić sobie jako taśmę
produkcyjną z~15~stanowiskami. W~każdym takcie algorytmu:
\begin{enumerate}[label=\arabic*.]
    \item Wszystko przesuwa się o~jedno miejsce w~lewo — bit ze stanowiska~0
          trafia na~1, z~1 na~2, ..., z~13 na~14. Bit ze stanowiska~14
          \emph{wylatuje} z~taśmy.
    \item Na zwolnione stanowisko~0 wstawiamy~0.
    \item Sprawdzamy, czy wylatujący bit~14 jest \emph{różny} (XOR) od bitu
          wejściowego \texttt{NXTBIT}. Jeśli tak — wykonujemy XOR całego
          rejestru z~\texttt{0x4599} (to jest właśnie „sprzężenie zwrotne"
          z~nazwy LFSR).
\end{enumerate}

Po przetworzeniu wszystkich 118~bitów wejścia zawartość rejestru to gotowa
suma kontrolna CRC-15.

Kluczowe szczegóły w~C++:
\begin{itemize}
    \item Rejestr to zmienna \texttt{uint16\_t} (16-bitowa liczba bez znaku).
          Używamy tylko 15~dolnych bitów — po każdym przesunięciu maskujemy:
          \texttt{rejestr \&= 0x7FFF} (zeruje bit~15, gdyby się pojawił po
          przesunięciu).
    \item Wejście wczytujemy jako znaki \texttt{'0'}/\texttt{'1'}.
          Szybka konwersja: \texttt{znak \& 1} daje wartość bitu — znaki
          \texttt{'0'} (kod ASCII 48) i~\texttt{'1'} (kod ASCII 49) różnią
          się tylko ostatnim bitem.
\end{itemize}

% ─────────────────────────────────────────────────────────────────────────────
\subsection*{B.3 Strategie optymalizacji algorytmu}

\subsubsection*{B.3.1 Eliminacja instrukcji warunkowej (kod bez skoków)}

Wewnętrzna pętla algorytmu zawiera \texttt{if}: gdy \texttt{CRCNXT}~=~1,
wykonaj XOR z~wielomianem; gdy~0, pomiń. Nowoczesne procesory działają
potokowo — zaczynają wykonywać następną instrukcję, zanim skończyły poprzednią.
Gdy napotykają \texttt{if}, muszą \emph{zgadnąć} który wariant będzie
potrzebny i~zacząć go wykonywać z~wyprzedzeniem. Przy losowych danych wejściowych
zgadują 50/50 — połowa zgadnięć jest błędna. Każdy błąd kosztuje 5–15~cykli
procesora na „cofnięcie" pracy i~ponowne wykonanie właściwej gałęzi.
Zjawisko to nosi nazwę \textit{branch misprediction} (błędne przewidywanie
gałęzi).

Rozwiązanie: zastąpić \texttt{if} operacją arytmetyczną dającą ten sam wynik
bez żadnego skoku:
\begin{itemize}
    \item \textbf{Mnożenie:} \texttt{CRCNXT * 0x4599} — daje \texttt{0x4599}
          gdy \texttt{CRCNXT}~=~1 i~\texttt{0} gdy \texttt{CRCNXT}~=~0.
    \item \textbf{Maska bitowa:} negacja \texttt{-CRCNXT} — daje
          \texttt{0xFFFF} gdy \texttt{CRCNXT}~=~1 i~\texttt{0x0000} gdy~=~0;
          XOR z~tą maską i~wielomianem zastępuje \texttt{if}.
\end{itemize}
Przy fladze \texttt{-O3} kompilator często sam wykrywa ten wzorzec i~generuje
kod bezskokowy.

\subsubsection*{B.3.2 Tablica przeglądowa (zamieniamy 8~kroków w~1)}

Algorytm w~wersji podstawowej przetwarza wejście \textbf{bit po bicie} —
96~iteracji pętli dla 96~bitów. Można go przyspieszyć 8-krotnie, przetwarzając
\textbf{bajt po bajcie} — 12~iteracji dla 96~bitów.

Idea: przed pętlą główną, raz na zawsze, oblicz co stanie się z~rejestrem po
przetworzeniu \emph{każdego} możliwego bajtu. Jest ich tylko $2^8 = 256$.
Zapisz wyniki w~tablicy 256~elementów. W~pętli głównej zamiast przetwarzać
8~bitów osobno — jedno spojrzenie do tablicy daje gotowy wynik.

Tablica ma rozmiar $256 \times 2$\,B = 512\,B, co mieści się w~całości
w~pamięci podręcznej L1 procesora (32–64\,kB). Po kilku pierwszych iteracjach
cała tablica jest już w~L1 — każdy odczyt trwa 1–4~cykli zamiast 100+~cykli
odczytu z~RAM.

Dzięki słowu kluczowemu \texttt{constexpr} tablica jest wypełniana podczas
\emph{kompilacji}, a~nie podczas uruchamiania programu — zero kosztu
inicjalizacji w~czasie uruchomienia.

\subsubsection*{B.3.3 Flagi kompilatora}

Kompilacja z~\texttt{-O3 -march=native} zleca kompilatorowi agresywną
optymalizację:
\begin{itemize}
    \item \textbf{Rozwijanie pętli} (\textit{loop unrolling}) — zamiast
          pętli z~12~lub~96 krokami kompilator generuje osobne instrukcje
          dla każdej iteracji, eliminując narzut skoków powrotnych.
    \item \textbf{Wbudowanie funkcji} (\textit{inlining}) — kompilator wkleja
          ciało funkcji CRC w~miejsce wywołania, eliminując narzut wywołania.
    \item \textbf{Instrukcje SIMD} — na Apple Silicon (NEON) i~x86-64 (AVX2)
          specjalne instrukcje operują na wielu danych jednocześnie.
\end{itemize}

% ─────────────────────────────────────────────────────────────────────────────
\subsection*{B.4 Pomiar czasu i~integralność benchmarku}

\subsubsection*{B.4.1 Mierzenie czasu}

Do pomiaru czasu w~C++ służy \texttt{std::chrono::high\_resolution\_clock}
z~nagłówka \texttt{<chrono>}. Zegar ten ma rozdzielczość rzędu nanosekund
na macOS i~Linux. Wzorzec użycia: zapisz czas startowy przed pętlą, czas
końcowy po pętli, odejmij i~przelicz na żądaną jednostkę (\texttt{milliseconds},
\texttt{nanoseconds} itp.) przez \texttt{duration\_cast}.
Dla wiarygodnych wyników powtórzeń $n$ powinno być co najmniej $10^6$ —
przy mniejszych wartościach zmierzony czas jest zdominowany przez szum
samego pomiaru zegara, a~nie przez rzeczywisty czas algorytmu.

\subsubsection*{B.4.2 Pułapka: kompilator może usunąć całą pętlę benchmarku}

Przy \texttt{-O3} kompilator analizuje kod i~może stwierdzić: „wywoływana
jest $n$~razy ta sama funkcja z~tym samym wejściem — wynik się nie zmienia,
więc obliczę go raz i~zastąpię całą pętlę jednym wywołaniem". Benchmark
pokazałby wtedy czas jednej iteracji niezależnie od wartości $n$ — zupełnie
błędny pomiar.

Dwa zabezpieczenia, które razem uniemożliwiają tę optymalizację:
\begin{itemize}
    \item \textbf{Zmienna \texttt{volatile}} — zapisanie wyniku CRC do zmiennej
          \texttt{volatile} jest przez standard C++ traktowane jako „widoczny
          efekt uboczny": kompilator nie może usunąć kodu, który do niej
          prowadzi, ani wynieść obliczeń poza pętlę.
    \item \textbf{\texttt{\_\_attribute\_\_((noinline))}} na funkcji CRC —
          blokuje wbudowanie (inlining) funkcji do miejsca wywołania.
          Kompilator nie „widzi" wtedy wnętrza funkcji i~nie może udowodnić,
          że wynik jest zawsze taki sam.
\end{itemize}
Oba mechanizmy stosujemy jednocześnie jako niezależne zabezpieczenia —
jedno może nie wystarczyć przy bardzo agresywnym optymalizatorze.

% ─────────────────────────────────────────────────────────────────────────────
\subsection*{B.5 Wyznaczanie $T_{\min}$ i~$T_{\max}$ — metoda obliczeniowa (zadanie 4.3)}

\subsubsection*{Parametry bazowe}

Przy prędkości 1\,Mb/s każdy bit trwa $T_b = 1\,\mu$s. Całkowity czas ramki
to liczba bitów pomnożona przez $T_b$.

Stała liczba bitów ramki (bez bitów stuffujących), z~tabeli w~A.2:

\begin{align*}
L_{\text{stały}} &= 1 + 32 + 6 + 64 + 16 + 2 + 7 + 3 = 131\ \text{bitów}
\end{align*}

Region objęty bit stuffingiem (patrz A.3):
\begin{align*}
N_{\text{fill}} &= 1 + 32 + 6 + 64 + 15 = 118\ \text{bitów}
\end{align*}

\subsubsection*{Najkrótszy czas — zero bitów stuffujących}

Żaden bit stuffujący nie zostanie dodany, gdy w~regionie 118~bitów nigdzie
nie pojawi się 5~identycznych bitów z~rzędu. Warunek ten spełniają dane
naprzemianne:
\begin{itemize}
    \item Dane (8~bajtów): \texttt{0x55 0x55 ...} (\texttt{0x55} =
          \texttt{01010101} — ciągłe przeskakiwanie 0 i~1)
    \item Identyfikator 29-bitowy: musi mieć naprzemienne bity
          (postaci \texttt{01010\ldots}), tak aby złożenie ID, pól Control
          i~wynikowego CRC nie tworzyło serii 5~identycznych na granicach pól
\end{itemize}

\begin{align*}
L_{\min} &= L_{\text{stały}} + 0 = \mathbf{131\ \text{bitów}} \\
T_{\min} &= 131 \cdot 1\,\mu\text{s} = \mathbf{131\ \mu\text{s}}
\end{align*}

Współczynnik efektywności — jaki procent czasu transmisji stanowią właściwe
dane (64~bity), a~jaki to nagłówki, sumy kontrolne i~inne pola ramki:
\begin{equation}
\eta_{\max} = \frac{64}{131} \approx \mathbf{48{,}9\%}
\end{equation}

\subsubsection*{Najdłuższy czas — maksimum bitów stuffujących (23)}

Maksymalne stuffowanie pojawia się, gdy co~5 bitów w~regionie wymusza
wstawienie bitu — wzorzec \texttt{11111\,00000\,11111\,00000}\ldots
Dla 118~bitów wejścia: $\lfloor 118 / 5 \rfloor = 23$~bity stuffujące.

Przykład danych generujących dużo stuffowania: same zera w~polu danych
\texttt{0x00} $\times$ 8 to 64~bity~\texttt{0} z~rzędu, co samo w~sobie daje
$\lfloor 64/5 \rfloor = 12$~bitów stuffujących. Odpowiedni dobór ID i~wynikowa
suma kontrolna CRC mogą łącznie dać pełne~23.

\textbf{Uwaga:} wartość CRC jest obliczana automatycznie z~ID i~danych —
nie można jej wybrać dowolnie. Znalezienie kombinacji dającej dokładnie 23~bity
stuffujące wymaga iteracyjnego sprawdzania kandydatów.

\begin{align*}
L_{\max} &= L_{\text{stały}} + 23 = 131 + 23 = \mathbf{154\ \text{bitów}} \\
T_{\max} &= 154 \cdot T_b = \mathbf{154\ \mu\text{s}}
\end{align*}

Współczynnik efektywności dla najdłuższej ramki:
\begin{equation}
\eta_{\min} = \frac{64}{154} \approx \mathbf{41{,}6\%}
\end{equation}

\subsubsection*{Podsumowanie zadania 4.3}

\begin{center}
\begin{tabular}{lrrrr}
\toprule
\textbf{Scenariusz} & \textbf{Bity stałe} & \textbf{Bity stuffuj.} & \textbf{Bity całk.} & $\bm{\eta}$ \\
\midrule
Najkrótsza ramka & 131 & 0  & 131 & 48,9\% \\
Najdłuższa ramka & 131 & 23 & 154 & 41,6\% \\
\bottomrule
\end{tabular}
\end{center}

Przy 1\,Mb/s: $T_{\min} = 131\,\mu$s, $T_{\max} = 154\,\mu$s.

\end{document}
